\chapter{Discussion} \label{chapter:ch6_discussion}

The introduction asked whether frequently changing collaborative-robot tasks can be programmed more directly by specifying spatial intent in the real workspace instead of relying on hand-guiding or editing motion commands and tool actions in a teach pendant or offline programming workflow. Within the \DIFdelbegin \DIFdel{scope of this prototype}\DIFdelend \DIFaddbegin \DIFadd{demonstrated thesis scope}\DIFaddend , the result is positive for \DIFdelbegin \DIFdel{selected }\DIFdelend \DIFaddbegin \DIFadd{prepared }\DIFaddend task-oriented workflows.

This chapter interprets that result rather than recapping the implementation. The thesis \DIFdelbegin \DIFdel{does not propose a }\DIFdelend \DIFaddbegin \DIFadd{develops a spatial authoring workflow for prepared use cases rather than a }\DIFaddend general robot programming environment. \DIFdelbegin \DIFdel{It shows that a selected use case }\DIFdelend \DIFaddbegin \DIFadd{Within that scope, it shows that the demonstrated use cases }\DIFaddend can combine workspace sensing, spatial input, task-specific interpretation, and robot execution into one authoring workflow, \DIFdelbegin \DIFdel{and it also makes clear where that result currently stops }\DIFdelend \DIFaddbegin \DIFadd{while also making explicit where that bounded result stops and where broader extension work begins}\DIFaddend .

\section{\DIFdelbegin \DIFdel{Prototype }\DIFdelend \DIFaddbegin \DIFadd{Thesis }\DIFaddend Result}

The main result is an end-to-end \DIFdelbegin \DIFdel{prototype }\DIFdelend \DIFaddbegin \DIFadd{implemented system }\DIFaddend for in-workspace spatial authoring on a collaborative robot. The \DIFdelbegin \DIFdel{implemented }\DIFdelend system combines a tracked 6DoF input device, stereo-vision-based scene sensing, task-specific interpretation, and robot execution behind a shared modular runtime with explicit interfaces. In practical terms, \DIFdelbegin \DIFdel{the prototype }\DIFdelend \DIFaddbegin \DIFadd{it }\DIFaddend covers the full path from setup and calibration through authoring, confirmation, visualization, execution, and persistence.

\DIFdelbegin \DIFdel{At prototype }\DIFdelend \DIFaddbegin \DIFadd{Within the thesis }\DIFaddend scope, this closes the three goals stated in the introduction. For \goalref{goal:1}, the thesis demonstrates that poses and trajectories can be authored directly in the workspace and turned into executable robot behavior through task-specific interpretation. For \goalref{goal:2}, it provides a modular system in which different use cases reuse the same sensing, orchestration, and execution infrastructure. For \goalref{goal:3}, it demonstrates that approach on a real robot with two representative task families, pick-and-place and seam welding, while \DIFdelbegin \DIFdel{keeping the evaluation qualitative}\DIFdelend \DIFaddbegin \DIFadd{evaluating the result qualitatively}\DIFaddend .

The \DIFdelbegin \DIFdel{important result is not the shared runtime alone, but }\DIFdelend \DIFaddbegin \DIFadd{shared runtime matters because of }\DIFaddend the common user workflow it enables. In both use cases, the operator follows the same overall process: choose the use case, provide the required spatial and non-spatial inputs, confirm the interpreted result, execute it, and save or load the authored program.

\DIFaddbegin \DIFadd{The strongest positive result appears after that pipeline has already been prepared. Chapter~\ref{chapter:ch5_user_workflow} makes clear that startup still requires module-graph construction, hardware activation, and camera calibration. The thesis does not aim to remove that preparation step. Its result is that, once the prepared pipeline is in place, repeated task instances can be authored through use-case-specific spatial input and confirmation instead of rebuilding low-level robot procedures for each new weld or pick target.
}

\DIFaddend That common path \DIFdelbegin \DIFdel{is }\DIFdelend \DIFaddbegin \DIFadd{also explains }\DIFaddend why the thesis did not produce only two isolated demonstrations. It produced a \DIFdelbegin \DIFdel{prototype }\DIFdelend \DIFaddbegin \DIFadd{system }\DIFaddend in which different task-oriented procedures \DIFdelbegin \DIFdel{can }\DIFdelend share the same acquisition, confirmation, execution, \DIFaddbegin \DIFadd{visualization, }\DIFaddend and persistence flow. \DIFdelbegin \DIFdel{The same structure also allows some components to change without forcing a redesign of the whole workflow, as long as the surrounding interfacesremain unchanged, although the current use case interfaces still limit how far that flexibility extends}\DIFdelend \DIFaddbegin \DIFadd{This shared structure is also a practical foundation for extension: the runtime already hosts two distinct use case modules on top of the same sensing and robot interfaces, and the same architecture can accommodate alternative sensing or authoring modules, or additional prepared use cases, without discarding the whole workflow}\DIFaddend .

That result is \DIFdelbegin \DIFdel{still intentionally bounded }\DIFdelend \DIFaddbegin \DIFadd{intentionally bounded by design}\DIFaddend . The system \DIFdelbegin \DIFdel{does not expose }\DIFdelend \DIFaddbegin \DIFadd{is built around prepared use cases rather than }\DIFaddend a full robot programming language, \DIFdelbegin \DIFdel{and it does not try to cover rich branching, loops, exception recovery, or complex sensor-driven logic. Its stronger result is narrower: some }\DIFdelend \DIFaddbegin \DIFadd{so it focuses on moving }\DIFaddend frequently changing tasks \DIFdelbegin \DIFdel{can be moved }\DIFdelend from repeated low-level program editing to re-authoring inside \DIFdelbegin \DIFdel{a use case that already defines }\DIFdelend \DIFaddbegin \DIFadd{use cases that already define }\DIFaddend most of the task structure.

\section{Pick-and-Place, Welding, and Use Case Scope}

The two demonstrated task families are not equivalent examples of the same success. They support different conclusions about where the approach is useful and what kind of task structure it fits best.

\DIFdelbegin \DIFdel{Pick-and-place is the clearest evidence that the prototype can combine authored intent with refreshed scene information at execution time. The authored command does not store only a fixed world-frame pick pose. It stores an object identity and }\DIFdelend \DIFaddbegin \DIFadd{Two concrete positive findings stand out. In the weld workflow, authoring is reduced to a scene scan, a traced pen trajectory along the intended seam, and confirmation of the interpreted result. The use case then derives the surrounding approach--weld--depart procedure from that trajectory and the detected seam geometry. In the pick-and-place workflow, the authored command stores the selected object identity together with }\DIFaddend a pick pose relative to that object, then \DIFaddbegin \DIFadd{refreshes the scene at execution time and }\DIFaddend reconstructs the world pick pose from \DIFdelbegin \DIFdel{a fresh scene estimate when the command is executed. This shows }\DIFdelend \DIFaddbegin \DIFadd{the stored object-relative transform. These are different strengths: weld shows direct authoring of the main task geometry, while pick-and-place shows scene-adaptive execution from authored spatial intent.
}

\DIFadd{Pick-and-place demonstrates }\DIFaddend that the system \DIFdelbegin \DIFdel{is not limited to replaying fixed authored motions . This result is meaningful, but only within the limits of the current scene representation and the specific pick-and-place task demonstrated in the prototype.
}\DIFdelend \DIFaddbegin \DIFadd{can combine authored spatial intent with refreshed scene information at execution time, rather than replaying fixed world-frame motions once taught.
}\DIFaddend 

Seam welding is the clearest practical fit for the present \DIFdelbegin \DIFdel{prototype. The authored trajectory expresses the main task intent }\DIFdelend \DIFaddbegin \DIFadd{system. Here, the changing part of the task is the weld geometry itself, and the operator can specify that geometry }\DIFaddend directly in the workspace \DIFdelbegin \DIFdel{, and }\DIFdelend \DIFaddbegin \DIFadd{while }\DIFaddend the use case \DIFdelbegin \DIFdel{can infer }\DIFdelend \DIFaddbegin \DIFadd{derives }\DIFaddend the surrounding robot procedure\DIFdelbegin \DIFdel{from that trajectory together with the detected scene geometry. In the demonstrated workflow, the user scans the scene, draws the weld, confirms the interpreted result, and then executes an approach--weld--depart sequence. That is a good match for }\DIFdelend \DIFaddbegin \DIFadd{. This is the closest match to }\DIFaddend the thesis question\DIFdelbegin \DIFdel{because it removes much of the }\DIFdelend \DIFaddbegin \DIFadd{, because it shifts the operator's work from rebuilding }\DIFaddend low-level motion \DIFdelbegin \DIFdel{specification while keeping the authored input concrete and easy to interpret}\DIFdelend \DIFaddbegin \DIFadd{structure to defining the task-relevant path itself}\DIFaddend .

The usefulness of this workflow depends on how often the weld changes. For stable repetitive production, existing programming workflows may already be sufficient once the task has already been programmed. The \DIFdelbegin \DIFdel{present approach becomes more }\DIFdelend \DIFaddbegin \DIFadd{thesis therefore suggests that the present approach is most }\DIFaddend useful when the overall task stays similar but the concrete weld changes often enough that repeated re-teaching becomes a noticeable burden, for example in small-batch production or frequently reworked jobs. In that setting, the human provides the task-specific intent in context and the robot performs the repeated physical execution.

More broadly, the thesis supports selected use cases with explicit task definitions\DIFdelbegin \DIFdel{, not general robot programming}\DIFdelend . In those cases, most of the task structure is prepared in advance by developing the use case module, and the operator defines the concrete spatial parameters of each task instance in the real workspace. A narrow use case can expose only the spatial inputs that actually vary and keep authoring direct. A broader use case can cover more variation, but it must also expose more parameters and more internal choices. The thesis does not resolve the best balance between those two directions, but it does show that both can be built on the same runtime structure as long as the task remains inside a prepared use case\DIFdelbegin \DIFdel{rather than a full programming environment}\DIFdelend \DIFaddbegin \DIFadd{. The broader question is therefore an architectural extension problem about how much variation a prepared use case should absorb before a more general programming model becomes necessary}\DIFaddend .

A workflow built around use cases could in principle be implemented without direct spatial input, for example by entering positions through a conventional interface. That would still simplify \DIFaddbegin \DIFadd{parts of }\DIFaddend the developed system, but it would leave the main spatial part of robot authoring in a more indirect form. For the task families considered here, poses and trajectories are not secondary parameters. They are a central part of what changes between task instances, so direct spatial input is part of the simplification rather than an optional addition.

\section{Applicability Boundary and Limitations}

The first major boundary is \DIFdelbegin \DIFdel{scene detection. The current prototype detects only }\DIFdelend \DIFaddbegin \DIFadd{the scene model used by the current system. It operates on }\DIFaddend predefined, tagged, box-shaped objects and builds the scene description from that narrow representation. This was sufficient for the \DIFdelbegin \DIFdel{feasibility-oriented tests used in the thesis , but it is also the main unresolved question for any broader practical use of the system}\DIFdelend \DIFaddbegin \DIFadd{demonstrated thesis tasks, but broader deployment would require richer scene understanding and corresponding validation}\DIFaddend . The thesis \DIFdelbegin \DIFdel{does not establish how reliable or general this scene detection would be outside the demonstrated setup. That limitation }\DIFdelend \DIFaddbegin \DIFadd{therefore does not target general object detection or broad-coverage scene modeling; instead, it evaluates the overall workflow in a deliberately bounded task-oriented setting. That boundary }\DIFaddend matters directly for both representative tasks, because the practical usefulness of the overall workflow depends strongly on the quality \DIFaddbegin \DIFadd{and breadth }\DIFaddend of the scene information it can obtain.

\DIFaddbegin \DIFadd{This boundary matters especially for welding. A broader direction considered during the design of the thesis was free-trajectory welding on existing workpieces with arbitrary seam shape. As discussed in }\autoref{subsec:ch2_analysis_scene_representation}\DIFadd{, supporting such scenarios would require a richer scene representation together with correspondingly richer downstream geometric processing. Under the current scene model, the demonstrated welding workflow is therefore naturally bounded to linear seams on the represented objects. The system architecture still leaves room for such extensions by separating sensing, scene interpretation, and use-case logic, but exploring richer scene pipelines would substantially broaden the technical scope, because scene acquisition, scene representation, and task-specific geometric reasoning each introduce significant design and implementation complexity.
}

\DIFaddend The next boundary is sensing and calibration robustness. Pen tracking depends on marker visibility, lighting conditions, motion blur, and corner detection quality, and the current pen implementation does not fuse the available IMU data. Camera calibration also assumes that the stereo rig remains rigid with respect to the robot flange. These are workable assumptions for the \DIFdelbegin \DIFdel{thesis prototype}\DIFdelend \DIFaddbegin \DIFadd{current implementation}\DIFaddend , but they limit how confidently the present sensing stack can be transferred to less controlled conditions.

Execution \DIFdelbegin \DIFdel{remains equally bounded}\DIFdelend \DIFaddbegin \DIFadd{scope is similarly focused}\DIFaddend . The thesis demonstrates a shared robot motion interface and shows that use cases can drive real robot execution through it\DIFdelbegin \DIFdel{, but generalized tool control is not implemented}\DIFdelend . In the evaluated use cases, gripper and welding actions are \DIFdelbegin \DIFdel{only }\DIFdelend approximated with fixed delays\DIFdelbegin \DIFdel{. No }\DIFdelend \DIFaddbegin \DIFadd{, while broader extensions such as generalized tool control, }\DIFaddend motion planning, collision checking, \DIFdelbegin \DIFdel{or simulated execution is provided. }\DIFdelend \DIFaddbegin \DIFadd{and simulated execution remain follow-up work for wider real-robot deployment rather than planned deliverables of the present thesis. In practice, because collaborative robot joints typically have limited ranges, repeated linear moves or arc motions could gradually drive the robot into a posture where the next move became infeasible, so safe }\type{moveJ} \DIFadd{transitions or a home-position reset were used as practical mitigations rather than as full motion planning.
}\DIFaddend 

The current use cases also depend on task-specific operating assumptions. The weld command assumes that the workpiece configuration stays unchanged after confirmation. The pick-and-place command assumes that the scene can be refreshed at execution time and that the expected object can still be matched by its stored identity.

UI maturity and portability are secondary limits. The frontend \DIFdelbegin \DIFdel{prototype }\DIFdelend supports setup, authoring, execution, and persistence in one workflow, but \DIFdelbegin \DIFdel{it still shows clear prototype-level behavior in practice}\DIFdelend \DIFaddbegin \DIFadd{some practical limitations remain}\DIFaddend . Visualization can require manual reload, saved setups may restore only partially if hardware modules cannot reactivate, and hardware support remains platform-limited: camera acquisition is currently implemented only on Windows, while BLE discovery is currently functional on Windows but not on Linux.

At the same time, the developed system's modular structure provides a practical path for further development, because individual parts of the system can be improved without requiring the whole \DIFdelbegin \DIFdel{prototype }\DIFdelend \DIFaddbegin \DIFadd{system }\DIFaddend to be rebuilt from the ground up. This does not remove the current limitations, and some improvements would still require changes at the interface boundary, but it does mean that the existing system remains a useful foundation for further work.

Taken together, these limits define the applicability boundary of the work. The thesis supports \DIFdelbegin \DIFdel{selected }\DIFdelend task-oriented workflows built around prepared use cases, not a general authoring method for arbitrary robot tasks. \DIFdelbegin \DIFdel{It does not yet provide the }\DIFdelend \DIFaddbegin \DIFadd{Broader deployment would benefit from richer }\DIFaddend scene understanding, \DIFdelbegin \DIFdel{execution safety, or tooling maturity required for broader deployment claims}\DIFdelend \DIFaddbegin \DIFadd{additional execution-safety layers, and further tooling maturation}\DIFaddend .

The evaluation boundary is equally explicit. The thesis demonstrates qualitative feasibility with \DIFdelbegin \DIFdel{limited physical }\DIFdelend \DIFaddbegin \DIFadd{targeted real-robot }\DIFaddend validation, and most development and testing took place in the Kassow REMII offline environment. It does not include a formal user study, a controlled comparison against teach pendant or offline programming workflows, or quantitative characterization of tracking accuracy, repeatability, or productivity. The result should therefore be read as \DIFdelbegin \DIFdel{a proof-of-concept for selected workflows, not as evidence of measured superiority or industrial readiness}\DIFdelend \DIFaddbegin \DIFadd{qualitative evidence that this authoring model can support prepared task-oriented workflows, while quantitative comparison and broader deployment evaluation remain follow-up work}\DIFaddend .
